\section*{Chapter \ref{ch:iv:what-each-interview-tests-by-role} --- What Each Interview Tests, by Role}

\begin{solution}{iq:iv:what-each-interview-tests-by-role:1}
The answer is 6. Let $a$ be the expected number of further tosses with no head pending and $b$ with one head
just seen: $a = 1 + \tfrac12 b + \tfrac12 a$ and $b = 1 + \tfrac12 a$, so $a = 6$. The trading interviewer
scores speed and a sense of spread: ``six; it is a waiting time, so it is skewed, and a run of fifteen tosses
would not surprise me'' (the variance is 22). The developer interviewer scores the method and a check: the two
equations, and either a ten-line simulation or the bounds (at least two tosses; at most a geometric wait of
pairs, $2 \times 4 = 8$, since disjoint pairs are HH with probability one quarter).
\iqlookfor{knowing that the same number is scored for speed and calibration in one room and for method and
verification in the other.}
\end{solution}

\begin{solution}{iq:iv:what-each-interview-tests-by-role:2}
(i) Whether the candidate turns a vague claim into a test with a null hypothesis: ``this signal predicts
returns; how would you check?''. (ii) Whether the candidate knows the ways a backtest lies (look-ahead,
survivorship, overfitting across many trials): ``of 200 variants tried, the top one shows a Sharpe ratio of
1.9; what do you conclude?''. (iii) Whether the candidate can reason about a model's assumptions and their
failure: ``your regression's coefficient flips sign when you add a variable; why?''. A trading interview asks
for a number now; a research interview asks how the candidate would know the number is right.
\iqlookfor{testing, overfitting and assumptions as the core of research, each tied to a concrete question.}
\end{solution}

\begin{solution}{iq:iv:what-each-interview-tests-by-role:3}
The offer probability per application is $\tfrac14 \times \tfrac12 \times \tfrac25 = \tfrac1{20}$. The number
of applications up to and including the first offer is geometric with mean 20. From twelve applications,
$1 - 0.95^{12} \approx 0.46$. A 90\% chance needs $n \ge \ln 0.1 / \ln 0.95 \approx 44.9$, so 45
applications. The independence assumption flatters the last figure: the same weakness fails several firms'
assessments, so improving the first stage (one in four) is worth more than applying more widely.
\iqlookfor{the product of stage rates, the geometric mean, the complement, and a word on independence.}
\end{solution}

\begin{solution}{iq:iv:what-each-interview-tests-by-role:4}
Correctness 3 or 4 (the final code is right; the bug was yours but you found it), complexity 2 or 3 (a working
but not optimal bound; say that you know an $O(n)$ answer may exist and what it would need), testing 4 (you
wrote the test that found the bug), communication 4. The single largest gain is on complexity: before coding,
ask what the lower bound is (every element must be read, so $O(n)$), and look for the structure (a hash map,
two pointers, a counting array) that reaches it. Five minutes spent there is worth more than the same five
minutes spent polishing.
\iqlookfor{self-assessment against the rubric's rows, and the habit of asking for the lower bound before
coding.}
\end{solution}

\begin{solution}{iq:iv:what-each-interview-tests-by-role:5}
Each interview is passed with probability $\Phi(0.5) \approx 0.691$. All four: $0.691^4 \approx 0.229$. At
least three: $4p^3(1-p) + p^4 \approx 0.637$. Passing all four with probability one half needs
$\Phi(\theta) = 0.5^{1/4} \approx 0.841$, so $\theta \approx 1.00$: a candidate one standard deviation of
noise above the bar is rejected half the time. A loop of vetoes rejects good candidates often and adds noise
rather than removing it; averaging scores (or requiring three of four) keeps the information of every
interview.
\iqlookfor{the binomial calculation, and the insight that a veto rule multiplies noise.}
\end{solution}

\begin{solution}{iq:iv:what-each-interview-tests-by-role:6}
A desk strategist's loop asks ``the desk wants a price for this new barrier variation by tomorrow; what do you
build tonight and what do you tell the trader it is wrong about?'' and ``hedge this book's vega with three
listed options''. A model validator's loop asks ``this local-volatility model reprices the calibration set
perfectly; what tests would still make you reject it?'' and ``design a benchmark model for an autocallable''.
The strategist is paid to deliver a usable number fast and to know its limits; the validator is paid to find
where a model fails and to be independent of the desk (One Quant Book~17, chapter~18).
\iqlookfor{the difference between building for the desk and challenging the desk's models.}
\end{solution}

\begin{solution}{iq:iv:what-each-interview-tests-by-role:7}
With ability spread 1 and noise 1, the average of $n$ independent scores has noise variance $1/n$, so its
correlation with ability is $1/\sqrt{1 + 1/n}$: 0.707, 0.816 and 0.894 for one, two and four interviewers.
If two interviewers share half their noise (a common component of variance $\tfrac12$ and private parts of
variance $\tfrac12$ each), the average's noise variance is $\tfrac12 + \tfrac14 = 0.75$ and the correlation
$1/\sqrt{1.75} \approx 0.756$: conferring first throws away most of the second interviewer. Rule: score
independently, record the scores, then discuss.
\iqlookfor{averaging independent noisy measurements, and why independence is the point.}
\end{solution}

\begin{solution}{iq:iv:what-each-interview-tests-by-role:8}
Test what the job does, in the order it matters. (i) Ten minutes: a validation question with time (``your
model's cross-validated score is 0.71 and live is 0.50; list the three most likely causes''), scored on
leakage, overlap and non-stationarity. (ii) Twenty-five minutes: code, turning a research function into a
production one: remove a look-ahead, add a test, make it streaming with bounded memory. (iii) Fifteen minutes:
design, serving a model with a latency budget, a fallback and monitoring for drift. (iv) Five minutes: the
candidate's questions. The rubric records, per part, correctness, whether the candidate tested or checked,
and whether the candidate named failure modes before being asked, each on anchored levels written before the
first candidate is seen.
\iqlookfor{a loop derived from the job's analysis, with anchored scores per dimension.}
\end{solution}
